\section{Signed compensation at logarithmic density}
\label{sec:signedjetcomparison}
The decisive improvement is to sum the complete first minor
correction with its signs before bounding the two-power
remainder. In particular a polynomial correction in the
selected size is not charged its absolute determinant mass.

\begin{lemma}[Signed occupancy moments at the parity root]
\label{lem:signedoccupancy}
Let \(T=J_N/N+B\) be real nonnegative doubly stochastic with
\(\|B\|\leq\rho\), \(0\leq\rho<1\), \(\max T\leq A/N\), \(A\geq2\).
For \(|I|=|J|=r\), let \(u=(\one_I,\one_J)\), and define
\[
 {\cal L}(f)=\sum_{r=1}^N(-1)^{r+1}r
       \sum_{|I|=|J|=r}\det T[I,J]^2 f(I,J).
\]
Put \(d(t)=\det(I-tBB^\top)\), \(d=d(1)>0\).
Uniformly over coordinates \(a,b\) of \(u\),
\[
 {\cal L}(1)=d,\quad
 \frac{{\cal L}(r(r-1))}{d}
                =\frac{4d'(1)+3d''(1)}d=O_\rho((1+A)^2),
\]
\[
 \frac{|{\cal L}(u_a)|}{d}\leq C_\rho(1+A)/N,\qquad
 \frac{|{\cal L}(u_a u_b)|}{d}
          \leq C_\rho(1+A)^2/N^2\quad(a\ne b).
\]
These are signed identities and bounds; no positivity of
individual occupancy moments is asserted.
\end{lemma}
\begin{proof}
Introduce separate row variables \(x_a\) on both parts,
let \(W=\diag(x)\), and put
\[
 F(t,x)=\det\left(I+\sqrt t\,W
                   \begin{pmatrix}0&T\\T^\top&0\end{pmatrix}\right),
 \qquad g(t,x)=-t\,\partial_tF(t,x).
\]
Cauchy--Binet shows that \(g(1,x)\) is precisely the signed
sum defining \({\cal L}\), with monomial \(x^u\).
At \(x=\one\),
\[
 F(t,\one)=(1-t)d(t),\qquad
 g(t,\one)=t d(t)-t(1-t)d'(t).
\]
Thus \(g(1,\one)=d\) and
\(g''(1,\one)=4d'(1)+3d''(1)\), giving the first identities.
The normalized derivatives of \(d\) through order two are
\(O_\rho((1+A)^2)\), by the convergent resolvent trace series.

For completeness we compute the coordinate derivatives at
the singular parity root. For \(t<1\), set
\[
 \Lambda(t)=\sqrt t\,P(I+\sqrt t\,P)^{-1},
       \qquad P=\begin{pmatrix}0&T\\T^\top&0\end{pmatrix}.
\]
With the unit constant and parity vectors \(c,s\) used above,
\[
 \Lambda(t)=\Gamma_{\sqrt t}
      +\frac{\sqrt t}{1+\sqrt t}cc^\top
      -\frac{\sqrt t}{1-\sqrt t}ss^\top .
\]
Write \(\tau=1-t\). Its expansion is
\[
 \Lambda(t)=-\frac2\tau ss^\top+H_0+\tau H_1+O(\tau^2),
 \qquad H_0=\Gamma+\tfrac12cc^\top+\tfrac32ss^\top .
\]
For clarity the coefficient of \(\tau\) is explicitly
\[
 H_1=-\tfrac12(\Gamma-\Gamma^2)
                  -\tfrac18cc^\top+\tfrac18ss^\top .
\]
Indeed \(\partial_s\Gamma_s|_{s=1}=\Gamma-\Gamma^2\).
The full block covariance satisfies the concrete bound
\[
 (1-\rho)|\Gamma_{ab}|\leq(2-\rho)A/N .
\]
To see this directly, \(\|\Omega\|\leq\|B\|\leq\rho\);
every row and column of \(\Omega\) has squared norm at most
\(A/N\). Writing \(R=(I+\Omega)^{-1}\), the identity
\(\Gamma=\Omega-\Omega R\Omega\) and
\(\|R\|\leq(1-\rho)^{-1}\) give the displayed estimate.
Also every row and column of \(\Gamma\) has norm at most
\(\sqrt{A/N}/(1-\rho)\). Consequently
\[
 |(\Gamma^2)_{ab}|\leq A/[N(1-\rho)^2] .
\]
Thus every entry of \(H_0,H_1\) is
\(O_\rho((1+A)/N)\), with the complete derivative and both
constant/parity contributions retained.

Let \(\partial_a\) here denote \(x_a\partial/\partial x_a\).
Multiaffinity in the rows gives, at \(x=\one\),
\[
 \partial_aF=F\Lambda_{aa},\qquad
 \partial_a\partial_bF
      =F(\Lambda_{aa}\Lambda_{bb}-\Lambda_{ab}^2)
                       \quad(a\ne b).
\]
The potential double pole in the second expression cancels
because the pole matrix \(ss^\top\) has rank one.
Define
\(q_{ab}(H)=s_a^2H_{bb}+s_b^2H_{aa}-2s_as_bH_{ab}\).
Taking the coefficient of \(\tau\), since \(g=-t\partial_tF\),
gives exactly
\[
 \frac{{\cal L}(u_a)}d=(H_0)_{aa}
                                  +2\frac{d'(1)}d s_a^2 ,
\]
and, for \(a\ne b\),
\[
 \frac{{\cal L}(u_a u_b)}d
 = (H_0)_{aa}(H_0)_{bb}-(H_0)_{ab}^2
       -2q_{ab}(H_1)+2\frac{d'(1)}d q_{ab}(H_0).
\]
Here \(s_a^2=1/(2N)\),
\(q_{ab}(H_i)=O_\rho((1+A)/N^2)\), and
\(d'(1)/d=O_\rho(1+A)\).
The stated coordinate bounds follow.
For a repeated coordinate use \(u_a^2=u_a\).
This evaluates every overlap type, including a row with
a column and two coordinates in the same part.
\end{proof}

\paragraph{Formalized finite occupancy input.}
The actual marked polynomials on the full vertex set \(V\) are
\[
 f_a(z)=\sum_{\substack{S\subseteq V\\a\in S}}\det P[S]z^{|S|},
 \qquad
 f_{ab}(z)=\sum_{\substack{S\subseteq V\\a,b\in S}}\det P[S]z^{|S|}.
\]
\texttt{HamiltonianCoordinateMinorPolynomial} proves these full
support expansions and their Euler readouts. The pair polynomial
includes the complete two-row inclusion--exclusion term, even when
the coordinates coincide. For distinct coordinates, writing \(L(z)=zP(I+zP)^{-1}\),
the actual row-replacement coefficient is
\(L_{aa}L_{bb}-L_{ab}L_{ba}\);
its off-diagonal subtraction is retained.

Put \(\Pi=ss^\top\) and, with \(z=1-\tau\), set
\[
 M(\tau)=I+(1-\tau)P+\Pi,\quad
 D(\tau)=\det M(\tau),\quad A(\tau)=I-M(\tau)^{-1}.
\]
The maintained proof derives that \(M(0)\) is invertible and
\(M(0)^{-1}=K+\Pi\) from the original stochastic and gap hypotheses.
For \(\tau\ne0,-1\), when \(M(\tau)\) is invertible, the complete
physical resolvent is exactly
\[
 L(1-\tau)=A(\tau)-\frac{\Pi}{\tau(1+\tau)} .
\]
Define, without suppressing either ordered cross term,
\[
 Q_{ab}(A)=\Pi_{aa}A_{bb}+\Pi_{bb}A_{aa}
                       -\Pi_{ab}A_{ba}-\Pi_{ba}A_{ab}.
\]
The rank-one identity
\(\Pi_{aa}\Pi_{bb}-\Pi_{ab}\Pi_{ba}=0\)
cancels the apparent squared pole. Consequently the actual marked
polynomials satisfy the exact cleared formulas
\[
 \begin{split}
 f_a(1-\tau)
   &=D(\tau)\left(\frac{\tau}{1+\tau}A_{aa}(\tau)
                         -\frac{\Pi_{aa}}{(1+\tau)^2}\right),\\
 f_{ab}(1-\tau)
   &=D(\tau)\left(\frac{\tau}{1+\tau}
         (A_{aa}A_{bb}-A_{ab}A_{ba})(\tau)
                         -\frac{Q_{ab}(A(\tau))}{(1+\tau)^2}\right)
                         \quad(a\ne b).
 \end{split}
\]
These identities, with precisely the displayed nonsingularity
conditions, are formalized in \texttt{HamiltonianPhysicalOccupancy},
using \texttt{HamiltonianParityRootRegularization},
\texttt{HamiltonianParityResolventSplit} and
\texttt{HamiltonianTwoRowDeterminant}.
The signed Euler readouts are sums over every principal support;
no independent coordinate model is introduced.
The regularized root jets now have maintained Lean proofs:
\[
 \begin{split}
 D(0)&=2d,\\
 D'(0)&=D(0)\bigl(\tfrac12-\operatorname{tr}\Gamma\bigr),\\
 A(0)&=\Gamma+\tfrac q4\one\one^\top,\\
 A'(0)&=-\Gamma+\Gamma^2-\tfrac q8\one\one^\top+\Pi,
       \qquad q=1/N.
 \end{split}
\]
The normalization is an exact constant-direction determinant
factorization. The derivative identity retains and cancels both
constant-direction cross products using stochastic balance.
The cleared formulas extend by their actual continuous limits;
the singular physical inverse is never evaluated at the root.

Precisely, define the full-principal-support functional
\[
 {\cal L}_{\rm pr}(f)=
   -\sum_{S\subseteq V}\frac{|S|}{2}\det P[S]\,f(S).
\]
The maintained polynomial Euler readouts and root derivatives prove
\[
 \frac{{\cal L}_{\rm pr}(u_a)}d
 =A_{aa}(0)+\bigl(\tfrac32+\operatorname{tr}\Gamma\bigr)\Pi_{aa},
\]
and, for \(a\ne b\),
\[
 \begin{split}
 \frac{{\cal L}_{\rm pr}(u_au_b)}d
 ={}&A_{aa}(0)A_{bb}(0)-A_{ab}(0)A_{ba}(0)
       -Q_{ab}(A'(0))\\
 &+\bigl(\tfrac32+\operatorname{tr}\Gamma\bigr)Q_{ab}(A(0)).
 \end{split}
\]
The Lean statements are the multiplied identities and require no
division by \(d\). Under the original nonnegative stochastic,
entry-cap and gap hypotheses, for every \(A\geq1\), they also prove
\[
 (1-\rho)|{\cal L}_{\rm pr}(u_a)|\leq5|d|Aq,\qquad
 (1-\rho)^2|{\cal L}_{\rm pr}(u_au_b)|
                         \leq50|d|A^2q^2\quad(a\ne b).
\]
In particular the covariance-square coefficient is independently
bounded by
\((1-\rho)^2|(\Gamma^2)_{ab}|\leq Aq\).
These are actual signed moments, rather than bounds on unsigned
minor mass. The maintained owners are
\texttt{HamiltonianParityRootJets},
\texttt{HamiltonianParityRootNormalization},
\texttt{HamiltonianClearedRootDerivative},
\texttt{HamiltonianRootOccupancyMoments},
\texttt{HamiltonianOccupancyEntryBounds},
\texttt{HamiltonianRootOccupancyBounds} and
\texttt{HamiltonianSignedOccupancyBounds}.
The arbitrary weighted balanced two-support readout and complete
joint first-order correction now also have maintained proofs.
The exact analytic size-derivative readout and the complete
remainder/tail and headline assembly remain unfinished.

\paragraph{Formalized complete signed main correction.}
For every real function on the finite vertex supports,
\texttt{HamiltonianSignedBipartiteFunctional} proves the actual identity
\[
 {\cal L}_{\rm pr}(f)=
 -\sum_{J\subseteq[1,N]}(-1)^{|J|}|J|
     \sum_{\substack{I\subseteq[1,N]\\|I|=|J|}}
          \det T[I,J]^2 f(I\sqcup J).
\]
The row and column supports are independent. Unbalanced determinants
are exactly zero, and the empty supports have their actual zero
rank coefficient. In particular the constant readout is \(d\).

Put \(h_a=-\Gamma_{aa}\) and \(v=(\tfrac12 I-K)h\).
The actual field satisfies
\((1-\rho)^2|v_a|\leq19A^2/N\).
\texttt{HamiltonianMinorFieldBudget} matches
\[
 \alpha(I,J)=r(r-1)/N+u^\top\Gamma u+2v^\top u
\]
to the precise first-order term of the maintained uniform minor
theorem. The finite signed functional is linear. The exact identities
\(u_a^2=u_a\) and
\[
 r(r-1)=\tfrac14\left(\sum_{\substack{a,b\in V\\a\ne b}}u_au_b
                                  -\sum_{a\in V}u_a\right)
\]
retain all overlaps and both orders of every off-diagonal pair.
With \(p=A/(1-\rho)\), the compiled full signed consumer gives
\[
 \left|{\cal L}(\alpha)\right|\leq900|d|p^3/N .
\]
Its proof controls the rank, diagonal, ordered off-diagonal and
full field contributions by
\[
 \frac{|d|}{N}
       \bigl(53p^2+20p^2+400p^3+380p^3\bigr)
             \leq\frac{900|d|p^3}{N},\qquad p\geq1.
\]
These bounds are taken after the actual signed coordinate
moments have been evaluated; no unsigned determinant mass enters
the main correction. The full-support finite identities are in
\texttt{HamiltonianSignedMinorFunctional}, and the actual stochastic
and balanced-support bounds are in
\texttt{HamiltonianJointSignedCorrection}. The separate complete
small-support residual now has a maintained quantitative proof,
as follows.

\paragraph{Formalized full small-support residual.}
Put \(R=\log(2N)^3\), \(r=|J|\), and
\[
 E_R(I,J)=\mathbf1_{r\leq R}
 \left[
  \left(\frac{\operatorname{per}T[I^c,J^c]}{\operatorname{per}T}\right)^2
                       -1-\alpha(I,J)\right].
\]
The actual retained error is the signed packet
\[
 \mathcal E_R=-\sum_{J\subseteq[1,N]}(-1)^r r
       \sum_{\substack{I\subseteq[1,N]\\|I|=r}}
                   \det T[I,J]^2 E_R(I,J).
\]
\texttt{HamiltonianAbsoluteMinorMass} evaluates its weighted mass
by the actual positive Gram spectrum and proves
\[
 \sum_J |J|\sum_{|I|=|J|}\det T[I,J]^2
                         \leq2A\det(I+B^\top B).
\]
The constant-direction factor is exactly \(2\).
\texttt{HamiltonianSignedMinorResidual} consumes the full uniform
two-power minor theorem. For every fixed \(0<\delta\leq1\)
and \(K\geq0\), one threshold covers every finite order,
every admissible kernel with \(A\leq(\log N)^K\), and all
independent deleted supports. It proves
\[
 |\mathcal E_R|\leq
 2A C_\delta^{\rm minor}
       \frac{\log(2N)^{12+24\lceil K\rceil}}{N^2}
                                  \det(I+B^\top B).
\]
The unsigned mass is used only for this residual. The previously
evaluated full first-order term retains its signed determinant
normalization.

\paragraph{Formalized reveal input to the large-tail estimate.}
For a finite family of injections separated by their complete row
histories, the sequential products
\[
 q_\sigma=\prod_{\text{revealed rows }i}
 \frac{M_{i,\sigma(i)}}{\sum_{j\ {\rm still\ available}}M_{ij}}
\]
have total mass at most \(1\).
\texttt{HamiltonianWeightedReveal} proves this with each actual
denominator and every branch retained. Define
\(w_\sigma=\prod_i M_{i,\sigma(i)}\) and let \(C_\sigma\)
be the sum of the logarithms of these denominators.
\texttt{HamiltonianWeightedRevealLog} proves, when the total
matching weight is positive,
\[
 \log\sum_\sigma w_\sigma
       \leq\frac{\sum_\sigma w_\sigma C_\sigma}
                      {\sum_\sigma w_\sigma}.
\]
It handles zero matching weights directly and derives the
inequality from \(\log x\leq x-1\).
The row-order average giving the global row-cap permanent bound,
its large-support complementary-permanent tail application, and
the final Hamiltonian comparison remain formalization obligations.

\begin{theorem}[Relative bipartite comparison at logarithmic density]
\label{thm:criticalrelative}
Fix \(0\leq\rho<1\) and \(C_0>0\).
Let \(T\) be real nonnegative doubly stochastic of order \(N\),
\[
 B=T-J_N/N,\quad \|B\|\leq\rho,\quad
 \max T\leq A/N,\quad 2\leq A\leq C_0\log(2N).
\]
Put
\[
 {\cal A}(T)=
        \frac{\det(I+BB^\top)}{\det(I-BB^\top)},\qquad
 P=\begin{pmatrix}0&T\\T^\top&0\end{pmatrix}.
\]
For some fixed finite \(M=M(\rho,C_0)\), uniformly in \(T\),
\[
 Z_{\rm H}(P)=\frac{(N!)^2}{2N\,N^{2N}}
 \left[1+O_{\rho,C_0}\left(
       \frac{\log(2N)^M}{N}
         +\frac{\log(2N)^M}{N^2}{\cal A}(T)\right)\right].
\]
In particular the full relative asymptotic holds whenever
\({\cal A}(T)\leq N^{2-\varepsilon}\) for fixed
\(\varepsilon>0\).
A sufficient explicit density condition is
\[
 A\leq a\log N,\qquad a\kappa_\rho<2,\qquad
 \kappa_\rho=\frac1{\rho^2}
                   \log\frac{1+\rho^2}{1-\rho^2};
\]
at \(\rho=0\) use the limiting value \(\kappa_0=2\).
\end{theorem}
\begin{proof}
Use the exact two-support Liu identity, normalized by
\((\operatorname{per}T)^2\).
On \(r\leq\log(2N)^3\), insert the complete expansion of
Theorem~\ref{thm:minorjet}.
The constant term is \(d=\det(I-BB^\top)\).
The \(r(r-1)/N\) term has normalized size \(O_\rho(A^2/N)\)
by Lemma~\ref{lem:signedoccupancy}.
The diagonal part of \(u^\top\Gamma u\) has normalized
size \(O_\rho(A^2/N)\), using all \(2N\) diagonal coordinates.
Its off-diagonal part has normalized size \(O_\rho(A^3/N)\):
there are at most \(4N^2\) ordered pairs, a coefficient
\(O_\rho(A/N)\), and a signed occupancy moment
\(O_\rho(A^2/N^2)\).
The field \(2v^\top u\) has the same normalized bound,
since \(\|v\|_\infty=O_\rho((1+A)^2/N)\).
These estimates are applied after the coupled signed moments
have been evaluated. The full main correction is consequently
\(O_\rho((1+A)^3/N)\), without an exponential determinant loss.

The moment identities cover all sizes.
Extending the truncated polynomial corrections to all \(r\)
costs only an exponentially small tail: their magnitude is
bounded by a polynomial in \(r,A,\log(2N)\), and the
Cauchy--Binet tail has \(r>\log(2N)^3\), \(A=O(\log N)\).
The same tail estimate with all complementary-permanent
ratios is
\(\exp(-c_{\rho,C_0}\log(2N)^3\log\log(2N))\),
as in the relative bipartite proof.

Only the second-order residual is now bounded by absolute
minor mass. Exactly, for \(Q=BB^\top\),
\[
 \sum_{r\geq1}r e_r(TT^\top)
 =\det(I+TT^\top)
       \operatorname{tr}\bigl(TT^\top(I+TT^\top)^{-1}\bigr)
 \leq2A\det(I+Q).
\]
After dividing by \(d\), its error is at most a fixed
logarithmic power times \(N^{-2}{\cal A}(T)\).
The determinant normalization remains explicit.

Thus
\[
 \frac{2N Z_{\rm H}(P)}{(\operatorname{per}T)^2}
 =d\left[1+O_{\rho,C_0}\left(
       \frac{\log(2N)^M}{N}
       +\frac{\log(2N)^M}{N^2}{\cal A}(T)\right)\right].
\]
The full permanent expansion has relative first correction
\(O_\rho(A^2/N)\).
Its squared reciprocal determinant cancels \(d\) exactly.
Multiplication gives the displayed relative formula; all product error terms are
absorbed by the displayed bound, since that permanent
relative correction tends to zero.

Finally the eigenvalues \(x\) of \(Q\) lie in \([0,\rho^2]\).
The function \(x^{-1}\log((1+x)/(1-x))\) is increasing there,
as its power series has nonnegative coefficients.
Also \(\operatorname{tr}Q\leq A-1\).
Consequently
\[
 \log{\cal A}(T)\leq\kappa_\rho\operatorname{tr}Q
                                      \leq\kappa_\rho A.
\]
The explicit sufficient condition leaves a positive power
between the two-power remainder and this amplification.
This proves the claimed relative range.
\end{proof}


\begin{corollary}[An inverse theorem for failure of relative universality]
\label{cor:relativeinverse}
Under Theorem~\ref{thm:criticalrelative}, fix
\(\eta>0\). If
\[
 \left|\frac{2N\,N^{2N}}{(N!)^2}Z_{\rm H}(P)-1\right|\geq\eta,
\]
then, for sufficiently large \(N\),
\[
 \log{\cal A}(T)\geq
            2\log N-O_{\rho,C_0,\eta}(\log\log(2N)),
\]
and consequently
\[
 \|T-J_N/N\|_F^2\geq
    \frac2{\kappa_\rho}\log N
              -O_{\rho,C_0,\eta}(\log\log(2N)).
\]
Thus a fixed relative counting defect requires a collective
nontrivial singular-value mass on the logarithmic scale.
A single fixed nontrivial eigenvalue cannot supply that
defect; all singular values are retained in \({\cal A}(T)\).
\end{corollary}
\begin{proof}
For large \(N\), the first error term
\(C_{\rho,C_0}\log(2N)^M/N\) is less than \(\eta/2\).
The second must therefore be at least \(\eta/2\).
It forces
\({\cal A}(T)\geq c_{\rho,C_0}\eta N^2/\log(2N)^M\).
Take logarithms, then use
\(\log{\cal A}(T)\leq\kappa_\rho\|B\|_F^2\).
This concerns a defect of the relative count, not just
Hamiltonian existence, which was already established in a
much wider diffuse range.
\end{proof}
